% Bagian: second-order-logic
% Bab: metatheory
% Subbagian: loewenheim-skolem

\documentclass[../../../include/open-logic-section]{subfiles}

\begin{document}

\olfileid[id]{sol}{met}{lst}

\olsection{Teorema L\"owenheim--Skolem Tidak Berlaku bagi Logika Orde Kedua}

\begin{explain}
Teorema L\"owenheim--Skolem (ke Bawah) menyatakan bahwa setiap himpunan
!!{sentence} yang mempunyai model tak hingga mempunyai model
!!{enumerable}. Teorema ini pun merupakan konsekuensi dari teorema
kelengkapan: pembuktian kelengkapan menghasilkan suatu model bagi setiap
himpunan !!{sentence} yang konsisten, dan model tersebut !!{enumerable}.
Terdapat pula Teorema L\"owenheim--Skolem ke Atas, yang menjamin bahwa jika
suatu himpunan !!{sentence} mempunyai model !!{denumerable}, maka himpunan
itu juga mempunyai model !!{nonenumerable}. Kedua teorema tersebut tidak
berlaku dalam logika orde kedua.
\end{explain}


\begin{thm}
\ollabel{thm:sol-no-ls} Teorema L\"owenheim--Skolem tidak berlaku bagi
logika orde kedua: terdapat !!{sentence} yang mempunyai model tak hingga,
tetapi tidak mempunyai model !!{enumerable}.
\end{thm}

\begin{proof}
Ingat bahwa
\[
\fn{Count} \ident \lexists[z][\lexists[u][\lforall[X][((X(z) \land
      \lforall[x][(X(x) \lif X(u(x)))]) \lif \lforall[x][X(x)])]]]
\]
benar dalam !!a{structure}~$\Struct{M}$ jika dan hanya jika $\Domain{M}$
!!{enumerable}, sehingga $\lnot\fn{Count}$ benar dalam~$\Struct{M}$ jika dan
hanya jika $\Domain{M}$ !!{nonenumerable}. Terdapat !!{structure} semacam
itu---ambil sembarang himpunan !!{nonenumerable} sebagai !!{domain}, misalnya
$\Pow{\Nat}$ atau~$\Real$. Jadi, $\lnot \fn{Count}$ mempunyai model tak
hingga, tetapi tidak mempunyai model !!{enumerable}.
\end{proof}

\begin{thm}
Terdapat !!{sentence} yang mempunyai model !!{denumerable}, tetapi tidak
mempunyai model !!{nonenumerable}.
\end{thm}

\begin{proof}
$\fn{Count} \land \fn{Inf}$ benar dalam $\Nat$, tetapi tidak benar dalam
  !!a{structure}~$\Struct{M}$ mana pun dengan $\Domain{M}$
!!{nonenumerable}.
\end{proof}


\end{document}
